% The sine, cosine and tangent of multidouble (src/trig.rs), generalized to n parts: an exact reduction by pi/2, Taylor
% series by Horner's scheme, and the bounds from Lemma 20 at every step, or from the published double-word bounds.
% \input after the sloppy section of expansion-error-bounds.tex, whose macros and labels it uses. Constants:
% trig_consts.py.

\section{Sine, cosine and tangent}\label{sec:trig}

The functions reduce $x$ by $\pi/2$ exactly, evaluate the Taylor series of $\sin r$ and $\cos r$ with the operations
of the earlier sections, and recombine. Throughout, $p = 53$, $n$ is even, $T_n = 2^{-53n/2}$, so that $T_n^2 =
u^n$, and $\bar a_i$ is that of \eqref{eq:abar} for $\theta = 2$.

\paragraph{Reduction} Let $F = 1216$, $P = \lfloor 2^F\pi/2 \rceil$ (an integer), $c = \RN(2/\pi)$ and $\tau =
\RN(0.7854)$, with $\lfloor\cdot\rceil$ the nearest integer, halves away from zero. RED takes a finite $n$-part $x$ with
$|x_0| \le 2^{52}$ and sets $k \gets \lfloor\RN(x_0c)\rceil$. If $k = 0$ it returns $q = 0$ and $r = x$. Otherwise
it forms the integer $N = 2^F\sum_i x_i - kP$ exactly. If $|\RN(2^{-F}N)| > \tau$, it sets $k_2 \gets
\lfloor\RN(\RN(2^{-F}N)\,c)\rceil$, $N \gets N - k_2P$ and $k \gets k + k_2$. It returns $q = k \bmod 4$ and the
greedy rounding $r$ of $2^{-F}N$ to $n$ parts: $r_0 = \RN(2^{-F}N)$ and $r_i = \RN(2^{-F}N - \sum_{l<i}r_l)$. In
multidouble $N$ is a two's-complement integer of 1280 bits. $P$ and $c$ come from Machin's formula in exact integer
arithmetic, checked against Gauss's formula (\texttt{scripts/trig\_tables.py}).

\paragraph{Evaluation} For $j \ge 0$ let $R_0 = 0.78541$ and $R_j = 2^{-j}(1 + 2^{-50})$ for $j \ge 1$. Bucket $j$ holds
the $r$ with $2^{-j-1} \le |r_0| < 2^{-j}$, and $[\frac12, \tau]$ for $j = 0$. $K_{\rm s}(j)$ is the least $K \ge 2$
with $R_j^{2K}/(2K+1)! \le 2^{-10}u^n(1 - R_j^2/6)$, and $K_{\rm c}(j)$ the least with $R_j^{2K}/(2K)! \le
2^{-10}u^n(1 - R_j^2/2)$. Let $s_k = (-1)^k/(2k+1)!$ and $c_k = (-1)^k/(2k)!$, and let $\hat s_k$ and $\hat c_k$ be
their greedy $n$-part roundings. For $|r_0| < T_n$, $\textrm{SINR}(r) = r$ and $\textrm{COSR}(r) = 1$. Otherwise,
with $z = \mathrm{MUL}(r, r)$, $K = K_{\rm s}(j)$ for SINR and $K_{\rm c}(j)$ for COSR,
\begin{gather*}
p_{K-1} = \hat s_{K-1}, \quad p_k = \mathrm{ADD}(\mathrm{MUL}(z, p_{k+1}), \hat s_k), \\
\textrm{SINR}(r) = 2^{-\sigma}\mathrm{ADD}\bigl(r', \mathrm{MUL}(r', \mathrm{MUL}(z, p_1))\bigr),
\end{gather*}
for $k = K-2, \dots, 1$, with $\sigma = j + 1$ and $r' = 2^\sigma r$, and COSR$(r) = \mathrm{ADDF}(\mathrm{MUL}(z,
p_1), 1)$ with $\hat c_k$ in place of $\hat s_k$. A power of two multiplies every part.

\paragraph{The functions} For $|x_0| < T_n$, SIN$_n(x) = x$, COS$_n(x) = 1$ and TAN$_n(x) = x$. Otherwise, with $(q, r)
= \mathrm{RED}(x)$, $Y_{\rm s} = \mathrm{SINR}(r)$ and $Y_{\rm c} = \mathrm{COSR}(r)$, SIN$_n$ returns $Y_{\rm s}$,
$Y_{\rm c}$, $-Y_{\rm s}$, $-Y_{\rm c}$ for $q = 0, 1, 2, 3$, COS$_n$ returns $Y_{\rm c}$, $-Y_{\rm s}$, $-Y_{\rm c}$,
$Y_{\rm s}$, and TAN$_n$ returns $2^{-3\sigma}\mathrm{DIV}(2^{3\sigma}Y_{\rm s}, Y_{\rm c})$ for even $q$ ($\sigma = 0$ for
$|r_0| < T_n$) and $-$DIV$(Y_{\rm c}, Y_{\rm s})$ for odd $q$. Unscaled, $r\,\mathrm{MUL}(z, p_1)$ has products below
$2^{-1022}$ for $|r|$ under about $2^{-189}$ at $n = 8$, and so have the products of $Y_{\rm c}$'s last parts, near
$2^{-2\sigma-318}$, with the last quotient terms for $|r|$ under about $2^{-81}$. At $n = 8$ these are multidouble's
\texttt{OD::sin}, \texttt{cos} and \texttt{tan}. Its \texttt{DD} versions run the same steps with DWTimesDW3 for MUL,
AccurateDWPlusDW for ADD, DWPlusFP for ADDF and DWDivDW2 for DIV~\cite{JMP17}, after a $\FastTwoSum$ of $r_0$ and $r_1$,
which makes $r$ a double-word number: the greedy $r$ can end in a tie, $|r_1| = \frac12\ulp(r_0)$ with $r_0$ odd.
The emulator of the earlier sections, running these steps, returns the bits of both on 600 inputs each, with ladders
at their top, results next to multiples of $\pi/2$ and arguments up to $2^{52}$. Controls that halve the sine series or
skip the correction of $k$ mismatch (\texttt{trig\_bits.py}).

\begin{lemma}\label{lem:red}
Assume~\ref{as:fp}, let $x$ be finite with $|x_0| \le 2^{52}$ satisfying $H(2)$, and let $(q, r) = \mathrm{RED}(x)$.
Then $|k| < 0.64\cdot2^{52}$, $q \equiv k \pmod 4$, and with $r_{\rm e} = x - k\pi/2$ and $A = |k|\,2^{-1217}$:
\begin{enumerate}
\item[(a)] if $k = 0$, $r = x$ and $|r| \le 0.785399$;
\item[(b)] otherwise $|r - r_{\rm e}| \le u_*^n|r_{\rm e}| + 2A$, $|r| \le R_0$, and $|r_{i+1}| \le \frac12\ulp(r_i)$,
  so $r$ satisfies $H(1)$.
\end{enumerate}
In both cases $r$ in bucket $j$ has $|r| \le R_j$, and its parts satisfy $|r_i| \le \bar a_i|r|$.
\end{lemma}

\begin{proof}
Each $x_i$ is a multiple of $2^{-1074}$, so $2^F\sum_i x_i$ is an integer, and $N$ is computed exactly; $|N|$ stays
below $2^{F+54}$. Write $v = 2^{-F}N$, so $v = r_{\rm e} + k(\pi/2 - 2^{-F}P)$ and $|v - r_{\rm e}| \le A$. By $H(2)$,
$|x - x_0| \le 2u|x_0|/(1-2u)$, so with $|c - 2/\pi| \le 2^{-54}$, \eqref{eq:rn} and $|x_0| \le 2^{52}$,
$|2x/\pi - \RN(x_0c)| \le 0.32 + 0.25 + 0.64 < 1.21$. So $|v| < 1.71\cdot\pi/2 + A < 2.7$ before the
correction, and $|k| \le (1+u)2^{52}c + \frac52 < 0.64\cdot2^{52}$ after it. If $|\RN(v)| \le \tau$, then $|v| \le
\tau(1+u)$. Otherwise the three terms of the same estimate with $|v| < 2.7$ total at most $5u$, so $|2v/\pi - k_2| \le
\frac12 + 5u$ and the new $|v| \le \pi/4 + 2^{-50} < \tau$. In both cases $|v| \le \tau(1+u) < R_0$. The correction
needs $|x_0| > 2^{31}$, as below it the first estimate is below $2^{-20}$ and $|v| < \pi/4 + 1.5\cdot10^{-6} < \tau$.
Then the first $k$ exceeds $2^{30}$ and $|k_2| \le 2$, so the final $k$ is 0 only when the first one is. (a)
$\lfloor\RN(x_0c)\rceil = 0$ means $|\RN(x_0c)| < \frac12$, so $|x_0| < \frac\pi4(1+u)/(1-u) < 0.7853982$, and $|x| \le
|x_0|/(1-2u)$. (b) By \eqref{eq:rn}, each
remainder of the greedy rounding is at most $u_*$ times the one before, so $|r - v| \le u_*^n|v|$, and $|r_{i+1}| \le
\frac12\ulp(r_i)$; then $|r - r_{\rm e}| \le u_*^n(|r_{\rm e}| + A) + A$. The bucket bounds follow from $|r| \le
|r_0|/(1-2u)$, and the part bounds are \eqref{eq:abar}, since $H(1)$ implies $H(2)$.
\end{proof}

The evaluation is a program of the operations of Lemma~\ref{lem:g}, and its bounds are the $G_{\rm OP}$ of that lemma
applied in turn, with $r$ in units of $|r|$, $z$ in units of $r^2$, and the Horner sums absolute. For bucket
$j$ with $R = R_j$ and $K = K_{\rm s}(j)$, let $\sigma_k$ be the magnitudes of the parts of $\hat s_k$, $\delta_k =
|\hat s_k - s_k|$, $(\epsilon_z, \zeta) = G_{\rm MUL}(\bar a, \bar a;\ 1)$, $E_{K-1} = \delta_{K-1}$ and $\pi^{(K-1)} =
\sigma_{K-1}$, and for $k = K-2, \dots, 1$
\begin{align*}
(\epsilon_m, \mu) &= G_{\rm MUL}\bigl(R^2\zeta, \pi^{(k+1)};\ R^2(1+\epsilon_z)(|s_{k+1}| + E_{k+1})\bigr), \\
D_k &= R^2(1+\epsilon_z)E_{k+1} + R^2\epsilon_z|s_{k+1}| + \delta_k + \epsilon_m, \\
(\epsilon_a, \pi^{(k)}) &= G_{\rm ADD}\bigl(\mu, \sigma_k;\ |s_k| + D_k\bigr), \quad E_k = D_k + \epsilon_a.
\end{align*}
Then let
\begin{align*}
(\epsilon_t, \tau') &= G_{\rm MUL}\bigl(\zeta, \pi^{(1)};\ (1+\epsilon_z)(|s_1| + E_1)\bigr), \\
E_t &= (1+\epsilon_z)E_1 + \epsilon_z|s_1| + \epsilon_t, \\
(\epsilon_w, \omega) &= G_{\rm MUL}\bigl(\bar a, \tau';\ |s_1| + E_t\bigr), \\
D &= R^2(E_t + \epsilon_w) + R^{2K}/(2K+1)!, \\
(\epsilon_y, \psi_{\rm s}) &= G_{\rm ADD}\bigl(\bar a, R^2\omega;\ 1 + D\bigr), \\
E_{\rm S} &= D + \epsilon_y - u^n,
\end{align*}
where $\epsilon_y - u^n(1 + D)$ is the last ADD's error beyond $u^n|r + w|$, and let $\delta_{\rm S} = u^n + E_{\rm
S}/(1 - R^2/6)$. For COSR, with $K = K_{\rm c}(j)$, the same formulas with
$\hat c_k$ and $c_k$ in place of $\hat s_k$ and $s_k$, and the magnitudes of the parts of $\hat c_k$ in place of
$\sigma_k$, give $E_1$, $\pi^{(1)}$, $\tau'$ and $E_t$, and
\begin{align*}
D_{\rm c} &= R^2E_t + R^{2K}/(2K)!, \\
(\epsilon_y, \psi_{\rm c}) &= G_{\rm ADDF}\bigl(R^2\tau', 1;\ 1 + D_{\rm c}\bigr), \\
E_{\rm C} &= D_{\rm c} + \epsilon_y - u^n,
\end{align*}
where $\epsilon_y - u^n(1 + D_{\rm c})$ is the last ADDF's error beyond $u^n$ times its exact sum, and let
$\delta_{\rm C} = u^n + E_{\rm C}/\cos_- R$, with $\cos_- R = 1 - R^2/2 + R^4/24 - R^6/720 \le \cos R$. For the
bucket $|r_0| < T_n$, $\delta_{\rm S} = T^2/(6 - T^2)$ and $\delta_{\rm C} = T^2/(2\cos_- T)$ with $T = T_n(1 +
2^{-50})$.

\begin{lemma}\label{lem:sincosr}
Assume~\ref{as:fp} and (A3), and let $r$ be in bucket $j$ with $|r| \le R_j$ and $|r_i| \le
\bar a_i|r|$. Then $|\textrm{SINR}(r) - \sin r| \le \delta_{\rm S}|\sin r|$ and $|\textrm{COSR}(r) - \cos r| \le
\delta_{\rm C}\cos r$, and the parts of SINR$(r)$ and COSR$(r)$ are at most $\psi_{\rm s}|r|$ and $\psi_{\rm c}$.
\end{lemma}

\begin{proof}
Let $z_{\rm e} = r^2$ and $p^{\rm e}_k = \sum_{l=k}^{K-1}s_lz_{\rm e}^{\,l-k}$. Since $z_{\rm e} \le R^2 < 1$, the terms
of $p^{\rm e}_k$ alternate in sign and decrease, so $|p^{\rm e}_k| \le |s_k|$. Lemma~\ref{lem:g} gives $|z - z_{\rm
e}| \le \epsilon_zr^2$, part bounds $\zeta r^2$, and $|z| \le (1+\epsilon_z)R^2$. By induction from $k = K-1$,
$|p_k - p^{\rm e}_k| \le E_k$ and $p_k$ has part bounds $\pi^{(k)}$: with $m = \mathrm{MUL}(z, p_{k+1})$,
\[
p_k - p^{\rm e}_k = z(p_{k+1} - p^{\rm e}_{k+1}) + (z - z_{\rm e})p^{\rm e}_{k+1} + (\hat s_k - s_k) + \hat\epsilon_m
+ \hat\epsilon_a,
\]
where $\hat\epsilon_m = m - zp_{k+1}$ and $\hat\epsilon_a = p_k - (m + \hat s_k)$ are bounded by Lemma~\ref{lem:g},
whose third arguments bound $|zp_{k+1}|$ and $|m + \hat s_k|$. The same expansion gives $|t - z_{\rm e}p^{\rm e}_1|
\le E_tr^2$ for $t = \mathrm{MUL}(z, p_1)$, and $|w - rz_{\rm e}p^{\rm e}_1| \le (E_t + \epsilon_w)|r|^3 \le
R^2(E_t + \epsilon_w)|r|$ for $w = \mathrm{MUL}(r, t)$, whose part bounds $\omega|r|^3$ are at most $R^2\omega|r|$.
The series' remainder after $K$ terms is at most its first omitted term, so $|r + rz_{\rm e}p^{\rm e}_1 - \sin r| \le
|r|R^{2K}/(2K+1)!$, and $|r + w - \sin r| \le D|r|$. Lemma~\ref{lem:g} bounds the error of the last ADD by $u^n|r +
w|$ plus a term that does not depend on its third argument, which is $\epsilon_y - u^n(1+D)$ in units of $|r|$, and
$|r + w| \le |\sin r| + D|r|$, so $|\textrm{SINR}(r) - \sin r| \le u^n|\sin r| + E_{\rm S}|r|$; and $|r| \le |\sin r|/(1 - R^2/6)$. SINR's last two steps run on $r' = 2^\sigma r$ and its result is scaled back: under
(A2) both scalings are exact, so in units of $2^\sigma|r|$ the bounds are those above. COSR is the same
with $|c_1| = \frac12$ and the absolute units of ADDF, and $\cos r \ge \cos_- R$. For $|r_0| < T_n$, $|\sin r - r| \le
|r|^3/6$ and $|1 - \cos r| \le r^2/2$.
\end{proof}

For the quotient, let $\theta_j \ge 1$ be the least value, rounded up, with $\psi_{{\rm s},i} \le (\theta_ju)^i\lambda_{\rm s}$
and $\psi_{{\rm c},i} \le (\theta_ju)^i\lambda_{\rm c}$ for $i \ge 1$, where $\lambda_{\rm s} = (1 - R^2/6)(1 -
\delta_{\rm S}) - \sum_{i \ge 1}\psi_{{\rm s},i}$ and $\lambda_{\rm c} = \cos_-R\,(1 - \delta_{\rm C}) - \sum_{i \ge
1}\psi_{{\rm c},i}$ bound the leading parts from below ($\theta_j = 2$ for $|r_0| < T_n$), and let $\delta_{\rm D} =
u^n + K_nu^{n+1}$ with $K_n$ that of Theorem~\ref{thm:t6} at $\theta = \theta_j$, which needs $\theta_ju \le 1/4$.

\begin{theorem}\label{thm:trig}
Assume~\ref{as:fp} with $p = 53$ and (A3), $n$ even, and let $x$ be finite with $|x_0| \le 2^{52}$
satisfying $H(2)$, $f \in \{\sin, \cos\}$, $A = |k|\,2^{-1217}$ and $\epsilon = 1 + 2u^n$. For each bucket $j$, let
\begin{align*}
C_j &= \epsilon\max\bigl(\delta_{\rm S} + \epsilon u^n,\ \delta_{\rm C} + \epsilon u^nR_j\tan_+ R_j\bigr), \\
C^{\rm T}_j &= \epsilon\Bigl(\max\Bigl(\frac{(1+\delta_{\rm S})(1+\delta_{\rm D})}{1 - \delta_{\rm C}},
  \frac{(1+\delta_{\rm C})(1+\delta_{\rm D})}{1 - \delta_{\rm S}}\Bigr) \\
&\qquad - 1 + \epsilon u^n\frac{2R_j}{\sin_-2R_j}\Bigr),
\end{align*}
with $\sin_-$ and $\sin_+$ the Taylor sums of $\sin$ ending with the terms of orders seven and five, and $\tan_+ =
\sin_+/\cos_-$, let $C$ and $C^{\rm T}$ be the largest $C_j$ and $C^{\rm T}_j$ over the buckets, including $|r_0| < T_n$, and at least
$T^2/(2\cos_-T)$ and $\tan_+(T)/T - 1$, the errors of the returns for $|x_0| < T_n$. Then
\begin{gather*}
|\mathrm{SIN}_n(x) - \sin x| \le C|\sin x| + 5A, \\
|\mathrm{COS}_n(x) - \cos x| \le C|\cos x| + 5A, \\
|\mathrm{TAN}_n(x) - \tan x| \le \Bigl(C^{\rm T} + \frac{5A}{|\sin x\cos x|}\Bigr)|\tan x|.
\end{gather*}
At $n = 8$, $C \le 3.163167\,u^8$ and $C^{\rm T} \le 6.418351\,u^8$. With the double-word bounds of DWTimesDW3,
$(4u^2 + u^3/2)/(1+u)^2$ \cite{MR22}, AccurateDWPlusDW, $3u^2/(1-4u)$, DWPlusFP, $2u^2$, and DWDivDW2, $15u^2 +
56u^3$ \cite{JMP17}, in place of Lemma~\ref{lem:g} and Theorem~\ref{thm:t6}, the bounds hold for multidouble's
double-double functions of a double-word $x$ ($x_0 = \RN(x)$) with $C \le 7.841411\,u^2$ and $C^{\rm T} \le
28.414065\,u^2$. Since $|k| < 0.64\cdot2^{52}$, $5A < 0.8\cdot2^{-1163}$.
\end{theorem}

\begin{proof}
If $|x_0| < T_n$, $|x| < T_n(1-2u)^{-1} \le T$, and the returned $x$, $1$ and $x$ err by at most $|x|^3/6$, $x^2/2$ and $\tan|x|
- |x|$, within the stated $C$ and $C^{\rm T}$. Otherwise Lemma~\ref{lem:red} gives $r$ in bucket $j$ with the
hypotheses of Lemma~\ref{lem:sincosr}, and $f(x) = \pm\sin r_{\rm e}$ or $\pm\cos r_{\rm e}$ by the quadrant, as SIN$_n$
and COS$_n$ return. Let $\eta = |r - r_{\rm e}| \le u_*^n|r_{\rm e}| + 2A$ (zero for $k = 0$). For the sine,
$|\sin r - \sin r_{\rm e}| \le \eta\cos\xi$ with $\xi$ between $r$ and $r_{\rm e}$, and $\cos\xi \le \cos|r_{\rm e}| +
\eta$; with $|r_{\rm e}|\cos|r_{\rm e}| \le \sin|r_{\rm e}|$ and $|r_{\rm e}| \le 1.12\sin|r_{\rm e}|$ this is at most
$u^n(1+u^n)|\sin r_{\rm e}| + 4A$. With Lemma~\ref{lem:sincosr} and $|\sin r| \le (1 + 1.2u^n)|\sin r_{\rm e}| + 2A$,
$|\mathrm{SINR}(r) - \sin r_{\rm e}| \le C_j|\sin r_{\rm e}| + 5A$, as $\delta_{\rm S} \le 1/4$. For the cosine,
$|\cos r - \cos r_{\rm e}| \le \eta(\sin|r_{\rm e}| + \eta) \le u^n(1+u^n)|r_{\rm e}|\tan|r_{\rm e}|\cos r_{\rm e} +
4A$, and the same steps give the second term of $C_j$. For the tangent, $|\tan r - \tan r_{\rm e}| = |\sin(r - r_{\rm
e})|/|\cos r\cos r_{\rm e}| \le \eta/|\cos r\cos r_{\rm e}|$, which relative to $|\tan r_{\rm e}|$ is at most
$\epsilon u^n\,2|r_{\rm e}|/\sin 2|r_{\rm e}| + 3A/|\sin r_{\rm e}\cos r_{\rm e}|$, and $2t/\sin 2t$ increases. With
$Y_{\rm s} = \sin r\,(1 + \hat\sigma)$, $Y_{\rm c} = \cos r\,(1 + \hat\gamma)$, $|\hat\sigma| \le \delta_{\rm S}$ and
$|\hat\gamma| \le \delta_{\rm C}$, the parts of $Y_{\rm s}$ and $Y_{\rm c}$ satisfy $H(\theta_j)$, so
Theorem~\ref{thm:t6} gives DIV$(Y_{\rm s}, Y_{\rm c}) = (Y_{\rm s}/Y_{\rm c})(1 + \hat d)$ with $|\hat d| \le
\delta_{\rm D}$, and $\tan r\,(1+\hat\sigma)(1+\hat d)/(1+\hat\gamma)$ is within the first term of $C^{\rm T}_j$.
The numerator's $2^{3\sigma}$ and the quotient's $2^{-3\sigma}$ are exact under (A2), and Theorem~\ref{thm:t6}'s bound
is relative, so the scaled quotient has it. For odd $q$ the quotient is $-$DIV$(Y_{\rm c}, Y_{\rm s})$, and $|\cot r -
\cot r_{\rm e}|/|\cot r_{\rm e}| = |\sin(r - r_{\rm e})|/(|\sin r|\cos r_{\rm e})$, with $|\sin r|$ where the even case
has $\cos r$. By (A2), $\RN(v)$, which is not zero as $P$ is odd and $|k| < 2^{142}$, is at least $2^{-1074} >
2^{90}A$, so $|\sin r| \ge (1 - \eta/|\sin r_{\rm e}|)|\sin r_{\rm e}|$ with $\eta/|\sin r_{\rm e}| \le 1.2u^n +
2^{-89}$. If $2A \le u^n|\sin r_{\rm e}|$ the factor is within $\epsilon$. Otherwise its excess over $\epsilon$, at
most $2^{-88}$, multiplies a term below $2A/|\sin r_{\rm e}|$, which the absolute term absorbs. The rest is as for even
$q$. For the double-double functions
the same recursion runs with each error bound $e_{\rm op}\bar Z$ in place of Lemma~\ref{lem:g}, and with the last
operation's $e_{\rm op}$ in place of $u^n$ in $E_{\rm S}$, $E_{\rm C}$, $\delta_{\rm S}$ and $\delta_{\rm C}$. Their
operands are double-word numbers, since every algorithm ends with a $\FastTwoSum$, $r$ passes through one, and the
constants are greedy.

The values are those of \texttt{trig\_consts.py}, in rational arithmetic rounded up to 200-bit dyadic numbers. It
evaluates one bucket per pair $(K_{\rm s}(j), K_{\rm c}(j))$, the one with the largest $R_j$: every bound above holds
for all $|r| \le R_j$, and the buckets with the same series lengths lie in its range. It checks that the series
lengths and constants it derives are those of multidouble's tables.
\end{proof}

\begin{remark}
The largest $C_j$ is that of the cosine in bucket 0, at $n = 8$: $0.785473\,u^8$ from the reduction, $u^8$ from the
last ADDF, and $1.377695\,u^8$ from the rest, mostly the product $t$, whose roundings of $z$, $p_1$ and $t$ each enter
with a weight near $R_0^2/2$, divided by $\cos R_0$. For $|r_0| < \frac12$ the bounds fall to $2.175822\,u^8$ and
$4.800526\,u^8$, and for $|r_0| < 2^{-5}$ to $2.000705\,u^8$ and $4.002771\,u^8$. Evaluated at other $n$, the bounds
are $3.168700\,u^2$ and $6.452138\,u^2$ at $n = 2$, $3.164154\,u^4$ and $6.426154\,u^4$ at $n = 4$, and
$3582.787311\,u^{12}$ and $4929.367335\,u^{12}$ at $n = 12$, where the second term of the multiplication's bound
dominates (Table~\ref{tab:main}). In the default run of multidouble's tests, random arguments, arguments near multiples
of $\pi/2$ with results down to $2^{-443}$, and arguments near $2^{52}$ give at most $0.48\,u^8$ for the octo-double
functions, the error $x^2/2$ of the early return of the cosine just below $T_8$, and below $0.01\,u^8$ for every other
kind of argument. The double-double ones give at most $1.82\,u^2$ for the sine and cosine and $4.07\,u^2$ for the
tangent. Whether the bounds are approached is not known.
\end{remark}

\begin{remark}\label{rem:a2}
The theorem assumes (A2) for the whole run. With the scalings, the operands of every operation have leading parts
between $2^{-427}$ and $2^{425}$: $r$ itself at least $T_8 = 2^{-212}$ in $z = \mathrm{MUL}(r, r)$, $r'$ and the sums
in SINR near 1, $z$ and $t$ at least $2^{-2\sigma-3}$, and the quotient near $2^{2\sigma}$, whose divisor $Y_{\rm c}$
has its last parts near $2^{-2\sigma-318}$. At $n = 8$ their products of order at most 8 then lie above $2^{-850}$,
and the two-products' exponents sum to more than $-800$, against the $-970$ that (A2) needs, so for parts as dense as
those of a greedy rounding no result of the run leaves the normal range, whatever $|r|$ is from $T_8$ to $R_0$. On the emulator, 32 random $r$ in each bucket from
$[2^{-212}, 2^{-211})$ to $[\frac12, \tau]$, greedy roundings and ladders at their top, give no result outside binary64
and no subnormal one in $z$, SINR, COSR or either quotient, where the unscaled SINR has such results from $|r| <
2^{-189}$ and the unscaled quotient from $|r| < 2^{-81}$ (\texttt{trig\_bits.py a2}). As for the other operations of
this paper, an input with a part far below the ladder can still make a product of two small parts underflow.
\end{remark}
